\SourcePage{468}{471}
\section{Continuous groups}

In addition to the finite point groups listed in §~93, there are also \emph{continuous point groups}, whose number of elements is infinite. These are the groups associated with \emph{axial} symmetry and with \emph{spherical} symmetry.

The simplest axial-symmetry group is $C_\infty$. It contains the rotations $C(\varphi)$ through any angle $\varphi$ about the symmetry axis (and is called the \emph{two-dimensional rotation group}). The groups $C_n$ approach this group in the limit $n\to\infty$. In the same manner, the limiting forms of $C_{nh}$, $C_{nv}$, $D_n$, and $D_{nh}$ are the continuous groups $C_{\infty h}$, $C_{\infty v}$, $D_\infty$, and $D_{\infty h}$, respectively.


A molecule has axial symmetry only when all its atoms lie on a single straight line. If it is not symmetric with respect to its midpoint, its point group is $C_{\infty v}$; besides rotations about the axis, this group also includes reflections $\sigma_v$ in every plane through the axis. If the molecule is symmetric with respect to its midpoint, its point group is $D_{\infty h}=C_{\infty v}\times C_i$. The groups $C_\infty$, $C_{\infty h}$, and $D_\infty$, on the other hand, cannot occur as molecular symmetry groups at all.

The group of complete spherical symmetry contains rotations through an arbitrary angle about every axis through the centre, together with reflections in every plane through that point. This group, which we shall denote by $K_h$, is the symmetry group of an isolated atom. It has as a subgroup the group $K$ of all rotations in space (called the three-dimensional rotation group, or simply the \emph{rotation group}). Adjoining a centre of symmetry to $K$ gives the group $K_h$:
\[
 K_h=K\times C_i.
\]



The elements of a continuous point group can be distinguished by one or several parameters ranging continuously over their possible values. For example, three Euler angles specifying a rotation of the coordinate system may serve as parameters of the rotation group.

The general properties of finite groups described in §~92, and the associated concepts (subgroups, conjugate elements, classes, and so forth), extend directly to continuous groups. Statements tied directly to the order of a group naturally cease to apply; one example is the assertion that the order of a subgroup divides the order of the group.

In $C_{\infty v}$ all symmetry planes are equivalent, and consequently all the reflections $\sigma_v$ form a single class containing a continuous

\SourcePage{469}{472}
set of elements. The symmetry axis is two-sided, so there is a continuous set of classes, each containing the two elements $C(\pm\varphi)$. The classes of $D_{\infty h}$ follow immediately from those of $C_{\infty v}$, since $D_{\infty h}=C_{\infty v}\times C_i$.

All axes in the rotation group $K$ are equivalent and two-sided. Its classes therefore consist of rotations through an angle of a specified absolute magnitude $|\varphi|$ about an arbitrary axis. The classes of $K_h$ follow directly from those of $K$.

The notion of representations, both reducible and irreducible, likewise extends directly to continuous groups. Every irreducible representation contains a continuous set of matrices, but the number of basis functions transformed into one another---the dimension of the representation---is finite. These functions can always be chosen so that the representation is unitary. A continuous group has infinitely many distinct irreducible representations, but they form a discrete sequence: they can be assigned successive numbers. The matrix elements and characters of these representations obey orthogonality relations generalizing the analogous relations for finite groups. In place of (94.9) we now have
\begin{equation}
 \int G_{ik}^{(\alpha)}G_{lm}^{(\beta)*}\,d\tau_G
 =\frac1{f_\alpha}\delta_{\alpha\beta}\delta_{il}\delta_{km}\int d\tau_G,
 \tag{98.1}
\end{equation}
and in place of (94.10),
\begin{equation}
 \int\chi^{(\alpha)}(G)\chi^{(\beta)}(G)^*\,d\tau_G
 =\delta_{\alpha\beta}\int d\tau_G.
 \tag{98.2}
\end{equation}
The integration in these formulas is the so-called invariant integration over the group. The integration element $d\tau_G$ is expressed in terms of the group parameters and their differentials in such a way that the action of any group transformation again produces the same integration element\footnote{The assertions made here about the properties of irreducible representations of continuous groups hold only when the integrals (98.1) and (98.2) converge; in particular, the ``group volume'' $\int d\tau_G$ must be finite. This condition is satisfied by continuous point groups. It is not satisfied, for example, by the so-called Lorentz group, which will be encountered in the relativistic theory.}. In the rotation group, for instance, one may take $d\tau_G=\sin\beta\,d\alpha\,d\beta\,d\gamma$, where $\alpha$, $\beta$, and $\gamma$ are the Euler angles specifying the rotation of the coordinate system (§~58); then $\int d\tau_G=8\pi^2$.

We have, in substance, already obtained the irreducible representations of the three-dimensional rotation group. We did so while finding the eigenvalues and eigenfunctions of the total angular momentum, although without employing the ter\SourcePageJoin{470}{473}minology of group theory. Apart from a constant factor, the operators for the components of angular momentum are the operators of infinitesimal rotations\footnote{In mathematical terminology these operators are called the \emph{generators} of the rotation group.}; the angular-momentum eigenvalues therefore characterize the behaviour of wavefunctions under rotations in space. For angular momentum $j$ there are $2j+1$ distinct eigenfunctions $\psi_{jm}$, distinguished by the values $m$ of its projection and belonging to a single energy level of degeneracy $2j+1$. A rotation of the coordinate system mixes these functions among themselves, and they thus furnish an irreducible representation of the rotation group. In group-theoretic terms, therefore, the numbers $j$ label the irreducible representations of the rotation group, with one representation of dimension $2j+1$ for every $j$. Since $j$ assumes integral and half-integral values, the dimensions $2j+1$ run through all positive integers $1,2,3,\ldots$


The basis functions of these representations were already examined in substance in §§~56 and 57 (and the representation matrices were found in §~58). A basis for the representation with a given $j$ consists of the $2j+1$ independent components of a symmetric spinor of rank $2j$ (equivalently, the set of $2j+1$ functions $\psi_{jm}$).

The irreducible representations of the rotation group belonging to half-integral $j$ have an essential special feature. On rotation through $2\pi$, their basis functions---the components of a spinor of odd rank---change sign. But a $2\pi$ rotation is the identity element of the group. We must therefore conclude that representations with half-integral $j$ are, as it is said, \emph{double-valued}: in such a representation each group element (a rotation about some axis through an angle $\varphi$, $0\leqslant\varphi\leqslant2\pi$) corresponds not to one matrix but to two matrices whose characters have opposite signs\footnote{It should be noted that double-valued representations of a group are not representations in the strict sense, since they are realized by multivalued basis functions; see also §~99.}.

As already noted, an isolated atom has the symmetry $K_h=K\times C_i$. From the standpoint of group theory, every atomic term is consequently associated with an irreducible representation of the rotation group $K$, which fixes the value of the atom's total angular

\SourcePage{471}{474}
momentum $J$, and with an irreducible representation of $C_i$, which fixes the parity of the state\footnote{The atomic Hamiltonian is also invariant under permutations of the electrons. In the non-relativistic approximation the spatial and spin wavefunctions separate, and one may speak of representations of the permutation group furnished by the spatial functions. Specifying an irreducible representation of the permutation group determines the total spin $S$ of the atom (see §~63). Once relativistic interactions are included, however, a wavefunction cannot be separated into spatial and spin parts. Symmetry under simultaneous permutations of the coordinates and spins of the particles supplies no characteristic of a term, since the Pauli principle permits only total wavefunctions that are antisymmetric in all the electrons. Correspondingly, when relativistic interactions are included, spin is not, strictly speaking, conserved; only the total angular momentum $J$ is conserved.}.



Placing an atom in an external electric field splits its energy levels. The number of distinct levels thereby produced, and the symmetry of their states, can be found by the method described in §~96. To do this, the reducible $(2J+1)$-dimensional representation of the symmetry group of the external field, furnished by the functions $\psi_{JM}$, must be decomposed into irreducible representations of that group. We therefore need the characters of the representation furnished by the functions $\psi_{JM}$. Since irreducible-representation characters have the same value for all elements of a class, it is enough to consider rotations about a single axis, chosen as the $z$ axis. Under a rotation through $\varphi$ about $z$, the wavefunctions $\psi_{JM}$ are, as we know, multiplied by $e^{iM\varphi}$, where $M$ is the projection of the angular momentum on that axis. The transformation matrix for the functions $\psi_{JM}$ is consequently diagonal, and its character is
\[
 \chi^{(J)}(\varphi)=\sum_{M=-J}^{J}e^{iM\varphi}
 =\frac{e^{i(J+1)\varphi}-e^{-iJ\varphi}}{e^{i\varphi}-1},
\]
or\footnote{To avoid misunderstanding, we emphasize that this formula uses a parametrization of group elements different from the Euler-angle parametrization: a transformation is specified by the direction of the rotation axis and by the angle $\varphi$ about it. With this parametrization it can be shown, for example, that the integration in (98.2) is to be performed with the element $2(1-\cos\varphi)\,d\varphi\,do$, where $do$ is the element of solid angle for the direction of the rotation axis.}
\begin{equation}
 \chi^{(J)}(\varphi)=\frac{\sin\!\left((J+\tfrac12)\varphi\right)}{\sin(\varphi/2)}.
 \tag{98.3}
\end{equation}
Under inversion $I$, on the other hand, all the functions $\psi_{JM}$ with different $M$ behave alike: they are multiplied by $+1$ or by $-1$ according as the atomic state is even or odd.



\SourcePage{472}{475}
Consequently, the character is
\begin{equation}
 \chi^{(J)}(I)=\pm(2J+1).
 \tag{98.4}
\end{equation}
Finally, the characters for reflection in a plane $\sigma$ and for a rotation-reflection through an angle $\varphi$ are calculated by writing these symmetry transformations in the form
\[
 \sigma=IC_2,\qquad S(\varphi)=IC(\pi+\varphi).
\]

We shall also consider the irreducible representations of the axial-symmetry group $C_{\infty v}$. This question was essentially answered already when we classified the electronic terms of a diatomic molecule, whose symmetry is precisely $C_{\infty v}$ when its two atoms are different. Two one-dimensional representations correspond to the terms $0^+$ and $0^-$ (the terms with $\Omega=0$): the identity representation $A_1$, and the representation $A_2$, in which the basis function is invariant under every rotation but changes sign under reflections in the planes $\sigma_v$. The doubly degenerate terms with $\Omega=1,2,\ldots$ correspond instead to two-dimensional representations, denoted by $E_1,E_2,\ldots$ Under a rotation through $\varphi$ about the axis their basis functions are multiplied by $e^{\pm i\Omega\varphi}$, while reflection in a plane $\sigma_v$ interchanges the two functions. The characters of all these representations are
\begin{equation}
\begin{array}{c|ccc}
 C_{\infty v}&E&2C(\varphi)&\infty\sigma_v\\ \hline
 A_1&1&1&1\\
 A_2&1&1&-1\\
 E_k&2&2\cos k\varphi&0
\end{array}
\tag{98.5}
\end{equation}
The irreducible representations of $D_{\infty h}=C_{\infty v}\times C_i$ follow directly from those of $C_{\infty v}$; they correspond to the classification of the terms of a diatomic molecule with identical nuclei.



If half-integral values are used for $\Omega$, the functions $e^{\pm i\Omega\varphi}$ realize double-valued irreducible representations of $C_{\infty v}$, corresponding to molecular terms with half-integral spin\footnote{Unlike the three-dimensional rotation group, here an appropriate choice of fractional values of $\Omega$ could produce not only single- and double-valued representations, but also triple-valued and still higher-valued ones. Physically possible eigenvalues of angular momentum, however, as the operator of an infinitesimal rotation in three dimensions, are determined by representations of the three-dimensional rotation group specifically. Thus triple-valued and higher-valued representations of the two-dimensional rotation group (and likewise of any finite symmetry group), although definable mathematically, have no physical meaning.}.
